% --- Archon physics formalization source begin ---
% archon:physics
% archon:covers PhyXMiniProblems/problem_phyx_mini_0585.lean
% archon:source-report reports/phyx_mini/problem_phyx_mini_0585.source.json

\chapter{Physics problem phyx\_mini\_0585}
\label{ch:PhyXMiniProblems_problem_phyx_mini_0585}

\paragraph{Problem source.}
\#\# Physical scenario

Figure shows a thin tube in which a finite potential trap has been set up where \textbackslash{}( V\_2 = 0\textbackslash{},\textbackslash{}text\{V\} \textbackslash{}). An electron is shown traveling rightward toward the trap, in a region with a voltage of \textbackslash{}( V\_1 = -9.00\textbackslash{},\textbackslash{}text\{V\} \textbackslash{}), where it has a kinetic energy of \textbackslash{}( 2.00\textbackslash{},\textbackslash{}text\{eV\} \textbackslash{}). When the electron enters the trap region, it can become trapped if it gets rid of enough energy by emitting a photon. The energy levels of the electron within the trap are \textbackslash{}( E\_1 = 1.0\textbackslash{},\textbackslash{}text\{eV\} \textbackslash{}), \textbackslash{}( E\_2 = 2.0\textbackslash{},\textbackslash{}text\{eV\} \textbackslash{}), and \textbackslash{}( E\_3 = 4.0\textbackslash{},\textbackslash{}text\{eV\} \textbackslash{}), and the nonquantized region begins at \textbackslash{}( E\_4 = 9.0\textbackslash{},\textbackslash{}text\{eV\} \textbackslash{}) as shown in the energy-level diagram of figure.

\#\# Question

What is the smallest energy (eV) such a photon can have?

\#\# Answer choices

- A: A: 2.0eV

- B: B: 4.0eV

- C: C: 7.00eV

- D: D: 5.0eV

\#\# Figure caption (auxiliary; use the image as primary evidence)

The image is an energy level diagram. It features a vertical axis labeled “Energy” with a zero point at the bottom. Four discrete energy levels are shown as horizontal lines, labeled \textbackslash{}(E\_1\textbackslash{}), \textbackslash{}(E\_2\textbackslash{}), \textbackslash{}(E\_3\textbackslash{}), and \textbackslash{}(E\_4\textbackslash{}) from bottom to top in light blue. Above \textbackslash{}(E\_4\textbackslash{}) is a shaded region labeled “Nonquantized,” indicating a continuous energy range. The nonquantized region is marked by a horizontal boundary at the top with a faded gradient suggesting increasing energy.

\#\# Dataset metadata

Quantum Phenomena; Physical Model Grounding Reasoning, Implicit Condition Reasoning

\paragraph{Recorded answer/context.}
C: C: 7.00eV

\paragraph{Figure/image path.}
/root/proposal\_for\_physic/science-mango/phyx\_mini\_run/phyx\_data/test\_image/585.png

\paragraph{Formalization target.}
create a compiling Lean file with sorry bodies at `PhyXMiniProblems/problem\_phyx\_mini\_0585.lean`. The Lean declarations must preserve the physical quantities, dimensions or dimensional roles, figure labels, governing-law hypotheses, and final relation expressed by this problem.
Use Mathlib/Physlib names found through LeanExplore where available. If a physics API is missing, introduce faithful local abstractions rather than scalar placeholder aliases.

\begin{theorem}[Physics formalization target]
\label{thm:physics:phyx_mini_0585:target}
\lean{PhyXMiniProblems.ProblemPhyXMini0585.smallestPhotonEnergy_is_seven_electronVolts}
\uses{def:physics:phyx-mini-0585:phyxminiproblems-problemphyxmini0585-energyinelectronvolts, def:physics:phyx-mini-0585:phyxminiproblems-problemphyxmini0585-electronfinitetrapsetup, def:physics:phyx-mini-0585:phyxminiproblems-problemphyxmini0585-matchesfinitepotentialtrapscenario, def:physics:phyx-mini-0585:phyxminiproblems-problemphyxmini0585-matchesproblemreadouts, def:physics:phyx-mini-0585:phyxminiproblems-problemphyxmini0585-matchessuppliedenergylevelfigure, def:physics:phyx-mini-0585:phyxminiproblems-problemphyxmini0585-hasphysicalfinitetrapparameters, def:physics:phyx-mini-0585:phyxminiproblems-problemphyxmini0585-satisfieselectronentryandcapturelaws, def:physics:phyx-mini-0585:phyxminiproblems-problemphyxmini0585-isleastcapturephotonenergy, def:physics:phyx-mini-0585:phyxminiproblems-problemphyxmini0585-answerchoice, def:physics:phyx-mini-0585:phyxminiproblems-problemphyxmini0585-answerchoice-iscorrectfor}
The assigned autoformalize agent should translate this physics problem into Lean declarations in the covered file, with theorem and lemma proof bodies written as `by sorry`.
\end{theorem}
\begin{proof}
This is an autoformalization task, not a proof task. Produce statements that can later be proved by the physics prover without weakening the physical model.
\end{proof}

% --- Archon declaration topology begin ---
\section{Declaration topology}
\begin{definition}[electric Potential Dimension]
  \label{def:physics:phyx-mini-0585:phyxminiproblems-problemphyxmini0585-electricpotentialdimension}
  \lean{PhyXMiniProblems.ProblemPhyXMini0585.electricPotentialDimension}
  Electric potential has dimension energy divided by electric charge
\end{definition}
\begin{proof}
  This is the declared model definition of electric Potential Dimension.
\end{proof}
\begin{definition}[Electric Potential Quantity]
  \label{def:physics:phyx-mini-0585:phyxminiproblems-problemphyxmini0585-electricpotentialquantity}
  \lean{PhyXMiniProblems.ProblemPhyXMini0585.ElectricPotentialQuantity}
  \uses{def:physics:phyx-mini-0585:phyxminiproblems-problemphyxmini0585-electricpotentialdimension}
  A signed, unit-independent physical electric potential
\end{definition}
\begin{proof}
  This is the declared model definition of Electric Potential Quantity.
\end{proof}
\begin{definition}[energy In Joules]
  \label{def:physics:phyx-mini-0585:phyxminiproblems-problemphyxmini0585-energyinjoules}
  \lean{PhyXMiniProblems.ProblemPhyXMini0585.energyInJoules}
  Read a physical energy in coherent-SI joules
\end{definition}
\begin{proof}
  This is the declared model definition of energy In Joules.
\end{proof}
\begin{definition}[energy In Electron Volts]
  \label{def:physics:phyx-mini-0585:phyxminiproblems-problemphyxmini0585-energyinelectronvolts}
  \lean{PhyXMiniProblems.ProblemPhyXMini0585.energyInElectronVolts}
  \uses{def:physics:phyx-mini-0585:phyxminiproblems-problemphyxmini0585-energyinjoules}
  Read a physical energy in electron volts. The calibration is Physlib's dimensionful constant DimEnergy.electronVolt, not a scalar alias for energy
\end{definition}
\begin{proof}
  This is the declared model definition of energy In Electron Volts.
\end{proof}
\begin{definition}[electric Potential In Volts]
  \label{def:physics:phyx-mini-0585:phyxminiproblems-problemphyxmini0585-electricpotentialinvolts}
  \lean{PhyXMiniProblems.ProblemPhyXMini0585.electricPotentialInVolts}
  \uses{def:physics:phyx-mini-0585:phyxminiproblems-problemphyxmini0585-electricpotentialquantity}
  Read a physical electric potential in coherent-SI volts
\end{definition}
\begin{proof}
  This is the declared model definition of electric Potential In Volts.
\end{proof}
\begin{definition}[Tube Region]
  \label{def:physics:phyx-mini-0585:phyxminiproblems-problemphyxmini0585-tuberegion}
  \lean{PhyXMiniProblems.ProblemPhyXMini0585.TubeRegion}
  The two constant-potential regions named in the prose
\end{definition}
\begin{proof}
  This is the declared model definition of Tube Region.
\end{proof}
\begin{definition}[Particle Species]
  \label{def:physics:phyx-mini-0585:phyxminiproblems-problemphyxmini0585-particlespecies}
  \lean{PhyXMiniProblems.ProblemPhyXMini0585.ParticleSpecies}
  Particle species relevant to the incident-beam description
\end{definition}
\begin{proof}
  This is the declared model definition of Particle Species.
\end{proof}
\begin{definition}[Tube Direction]
  \label{def:physics:phyx-mini-0585:phyxminiproblems-problemphyxmini0585-tubedirection}
  \lean{PhyXMiniProblems.ProblemPhyXMini0585.TubeDirection}
  Direction of motion along the thin tube
\end{definition}
\begin{proof}
  This is the declared model definition of Tube Direction.
\end{proof}
\begin{definition}[Capture Mechanism]
  \label{def:physics:phyx-mini-0585:phyxminiproblems-problemphyxmini0585-capturemechanism}
  \lean{PhyXMiniProblems.ProblemPhyXMini0585.CaptureMechanism}
  Physical mechanism by which the electron loses energy and is captured
\end{definition}
\begin{proof}
  This is the declared model definition of Capture Mechanism.
\end{proof}
\begin{definition}[Trap Energy Marker]
  \label{def:physics:phyx-mini-0585:phyxminiproblems-problemphyxmini0585-trapenergymarker}
  \lean{PhyXMiniProblems.ProblemPhyXMini0585.TrapEnergyMarker}
  The four labels printed beside the horizontal lines in the supplied diagram. E₄ is the continuum boundary marker rather than an additional bound state
\end{definition}
\begin{proof}
  This is the declared model definition of Trap Energy Marker.
\end{proof}
\begin{definition}[Bound Trap Level]
  \label{def:physics:phyx-mini-0585:phyxminiproblems-problemphyxmini0585-boundtraplevel}
  \lean{PhyXMiniProblems.ProblemPhyXMini0585.BoundTrapLevel}
  The three final levels at which the electron is genuinely trapped
\end{definition}
\begin{proof}
  This is the declared model definition of Bound Trap Level.
\end{proof}
\begin{definition}[to Energy Marker]
  \label{def:physics:phyx-mini-0585:phyxminiproblems-problemphyxmini0585-boundtraplevel-toenergymarker}
  \lean{PhyXMiniProblems.ProblemPhyXMini0585.BoundTrapLevel.toEnergyMarker}
  \uses{def:physics:phyx-mini-0585:phyxminiproblems-problemphyxmini0585-trapenergymarker, def:physics:phyx-mini-0585:phyxminiproblems-problemphyxmini0585-boundtraplevel}
  Regard a bound-state label as its marker in the four-line diagram
\end{definition}
\begin{proof}
  This is the declared model definition of to Energy Marker.
\end{proof}
\begin{definition}[Spectrum Kind]
  \label{def:physics:phyx-mini-0585:phyxminiproblems-problemphyxmini0585-spectrumkind}
  \lean{PhyXMiniProblems.ProblemPhyXMini0585.SpectrumKind}
  Spectral character of an energy range in the finite trap
\end{definition}
\begin{proof}
  This is the declared model definition of Spectrum Kind.
\end{proof}
\begin{definition}[Trap Energy Level Figure]
  \label{def:physics:phyx-mini-0585:phyxminiproblems-problemphyxmini0585-trapenergylevelfigure}
  \lean{PhyXMiniProblems.ProblemPhyXMini0585.TrapEnergyLevelFigure}
  \uses{def:physics:phyx-mini-0585:phyxminiproblems-problemphyxmini0585-trapenergymarker}
  Literal visual information transcribed from phyx\_data/test\_image/585.png. The image itself displays labels and geometry, while the numerical level readouts are supplied by the prose and live in MatchesProblemReadouts
\end{definition}
\begin{proof}
  This is the declared model definition of Trap Energy Level Figure.
\end{proof}
\begin{definition}[Electron Finite Trap Setup]
  \label{def:physics:phyx-mini-0585:phyxminiproblems-problemphyxmini0585-electronfinitetrapsetup}
  \lean{PhyXMiniProblems.ProblemPhyXMini0585.ElectronFiniteTrapSetup}
  \uses{def:physics:phyx-mini-0585:phyxminiproblems-problemphyxmini0585-electricpotentialquantity, def:physics:phyx-mini-0585:phyxminiproblems-problemphyxmini0585-tuberegion, def:physics:phyx-mini-0585:phyxminiproblems-problemphyxmini0585-particlespecies, def:physics:phyx-mini-0585:phyxminiproblems-problemphyxmini0585-tubedirection, def:physics:phyx-mini-0585:phyxminiproblems-problemphyxmini0585-capturemechanism, def:physics:phyx-mini-0585:phyxminiproblems-problemphyxmini0585-trapenergymarker, def:physics:phyx-mini-0585:phyxminiproblems-problemphyxmini0585-boundtraplevel, def:physics:phyx-mini-0585:phyxminiproblems-problemphyxmini0585-spectrumkind, def:physics:phyx-mini-0585:phyxminiproblems-problemphyxmini0585-trapenergylevelfigure}
  Independent physical observables in the potential-trap experiment. In particular, neither electronEnergyUponEnteringTrap nor any emitted photon energy is defined from a numerical answer; both are constrained only by the governing laws below
\end{definition}
\begin{proof}
  This is the declared model definition of Electron Finite Trap Setup.
\end{proof}
\begin{definition}[Matches Finite Potential Trap Scenario]
  \label{def:physics:phyx-mini-0585:phyxminiproblems-problemphyxmini0585-matchesfinitepotentialtrapscenario}
  \lean{PhyXMiniProblems.ProblemPhyXMini0585.MatchesFinitePotentialTrapScenario}
  \uses{def:physics:phyx-mini-0585:phyxminiproblems-problemphyxmini0585-boundtraplevel, def:physics:phyx-mini-0585:phyxminiproblems-problemphyxmini0585-boundtraplevel-toenergymarker, def:physics:phyx-mini-0585:phyxminiproblems-problemphyxmini0585-electronfinitetrapsetup}
  Qualitative apparatus and particle roles stated in the problem
\end{definition}
\begin{proof}
  This is the declared model definition of Matches Finite Potential Trap Scenario.
\end{proof}
\begin{definition}[Matches Problem Readouts]
  \label{def:physics:phyx-mini-0585:phyxminiproblems-problemphyxmini0585-matchesproblemreadouts}
  \lean{PhyXMiniProblems.ProblemPhyXMini0585.MatchesProblemReadouts}
  \uses{def:physics:phyx-mini-0585:phyxminiproblems-problemphyxmini0585-energyinelectronvolts, def:physics:phyx-mini-0585:phyxminiproblems-problemphyxmini0585-electricpotentialinvolts, def:physics:phyx-mini-0585:phyxminiproblems-problemphyxmini0585-electronfinitetrapsetup}
  Numerical readouts supplied in the prose. These are calibrated measurements of physical quantities and contain no emitted-photon energy or minimizer
\end{definition}
\begin{proof}
  This is the declared model definition of Matches Problem Readouts.
\end{proof}
\begin{definition}[Matches Supplied Energy Level Figure]
  \label{def:physics:phyx-mini-0585:phyxminiproblems-problemphyxmini0585-matchessuppliedenergylevelfigure}
  \lean{PhyXMiniProblems.ProblemPhyXMini0585.MatchesSuppliedEnergyLevelFigure}
  \uses{def:physics:phyx-mini-0585:phyxminiproblems-problemphyxmini0585-electronfinitetrapsetup}
  All unambiguous primary-image evidence: an energy axis and zero, four labeled horizontal light-blue lines, and a shaded nonquantized region beginning at the E₄ marker and fading upward
\end{definition}
\begin{proof}
  This is the declared model definition of Matches Supplied Energy Level Figure.
\end{proof}
\begin{definition}[Has Physical Finite Trap Parameters]
  \label{def:physics:phyx-mini-0585:phyxminiproblems-problemphyxmini0585-hasphysicalfinitetrapparameters}
  \lean{PhyXMiniProblems.ProblemPhyXMini0585.HasPhysicalFiniteTrapParameters}
  \uses{def:physics:phyx-mini-0585:phyxminiproblems-problemphyxmini0585-energyinelectronvolts, def:physics:phyx-mini-0585:phyxminiproblems-problemphyxmini0585-boundtraplevel, def:physics:phyx-mini-0585:phyxminiproblems-problemphyxmini0585-boundtraplevel-toenergymarker, def:physics:phyx-mini-0585:phyxminiproblems-problemphyxmini0585-electronfinitetrapsetup}
  Sign and ordering conditions for a physical finite trap. These conditions exclude negative photon energies but do not select a minimizing level or give the requested 7 eV value
\end{definition}
\begin{proof}
  This is the declared model definition of Has Physical Finite Trap Parameters.
\end{proof}
\begin{definition}[Satisfies Electron Entry And Capture Laws]
  \label{def:physics:phyx-mini-0585:phyxminiproblems-problemphyxmini0585-satisfieselectronentryandcapturelaws}
  \lean{PhyXMiniProblems.ProblemPhyXMini0585.SatisfiesElectronEntryAndCaptureLaws}
  \uses{def:physics:phyx-mini-0585:phyxminiproblems-problemphyxmini0585-energyinelectronvolts, def:physics:phyx-mini-0585:phyxminiproblems-problemphyxmini0585-electricpotentialinvolts, def:physics:phyx-mini-0585:phyxminiproblems-problemphyxmini0585-boundtraplevel, def:physics:phyx-mini-0585:phyxminiproblems-problemphyxmini0585-boundtraplevel-toenergymarker, def:physics:phyx-mini-0585:phyxminiproblems-problemphyxmini0585-electronfinitetrapsetup}
  The governing laws used by the calculation. For an electron, a rise of ΔV volts in electrostatic potential lowers its potential energy by ΔV electron volts, so the same amount is added to its energy upon entering the trap. Single-photon capture then conserves energy: the entry energy is the sum of the final bound energy and emitted photon energy. Neither law specializes the result to 11 eV, 7 eV, or E₃
\end{definition}
\begin{proof}
  This is the declared model definition of Satisfies Electron Entry And Capture Laws.
\end{proof}
\begin{definition}[Is Least Capture Photon Energy]
  \label{def:physics:phyx-mini-0585:phyxminiproblems-problemphyxmini0585-isleastcapturephotonenergy}
  \lean{PhyXMiniProblems.ProblemPhyXMini0585.IsLeastCapturePhotonEnergy}
  \uses{def:physics:phyx-mini-0585:phyxminiproblems-problemphyxmini0585-energyinelectronvolts, def:physics:phyx-mini-0585:phyxminiproblems-problemphyxmini0585-electronfinitetrapsetup}
  An energy is realizable and no larger than the photon emitted for capture at any of the three bound levels. This predicate does not define an energy by a closed-form answer; it compares the independent photon observables in the setup
\end{definition}
\begin{proof}
  This is the declared model definition of Is Least Capture Photon Energy.
\end{proof}
\begin{lemma}[electron Energy Upon Entering Trap is eleven electron Volts]
  \label{lem:physics:phyx-mini-0585:phyxminiproblems-problemphyxmini0585-electronenergyuponenteringtrap-is-eleven-electronvolts}
  \lean{PhyXMiniProblems.ProblemPhyXMini0585.electronEnergyUponEnteringTrap_is_eleven_electronVolts}
  \uses{def:physics:phyx-mini-0585:phyxminiproblems-problemphyxmini0585-energyinelectronvolts, def:physics:phyx-mini-0585:phyxminiproblems-problemphyxmini0585-electronfinitetrapsetup, def:physics:phyx-mini-0585:phyxminiproblems-problemphyxmini0585-matchesfinitepotentialtrapscenario, def:physics:phyx-mini-0585:phyxminiproblems-problemphyxmini0585-matchesproblemreadouts, def:physics:phyx-mini-0585:phyxminiproblems-problemphyxmini0585-satisfieselectronentryandcapturelaws}
  The electron reaches the trap with 11 eV available for capture
\end{lemma}
\begin{proof}
  The conclusion follows from the governing relations and figure readouts named in the statement of electron Energy Upon Entering Trap is eleven electron Volts.
\end{proof}
\begin{lemma}[photon Energy For Capture At E3 is seven electron Volts]
  \label{lem:physics:phyx-mini-0585:phyxminiproblems-problemphyxmini0585-photonenergyforcaptureate3-is-seven-electronvolts}
  \lean{PhyXMiniProblems.ProblemPhyXMini0585.photonEnergyForCaptureAtE3_is_seven_electronVolts}
  \uses{def:physics:phyx-mini-0585:phyxminiproblems-problemphyxmini0585-energyinelectronvolts, def:physics:phyx-mini-0585:phyxminiproblems-problemphyxmini0585-electronfinitetrapsetup, def:physics:phyx-mini-0585:phyxminiproblems-problemphyxmini0585-matchesfinitepotentialtrapscenario, def:physics:phyx-mini-0585:phyxminiproblems-problemphyxmini0585-matchesproblemreadouts, def:physics:phyx-mini-0585:phyxminiproblems-problemphyxmini0585-satisfieselectronentryandcapturelaws}
  Capture into the highest bound state E₃ emits a 7 eV photon
\end{lemma}
\begin{proof}
  The conclusion follows from the governing relations and figure readouts named in the statement of photon Energy For Capture At E3 is seven electron Volts.
\end{proof}
\begin{lemma}[capture At E3 has least photon Energy]
  \label{lem:physics:phyx-mini-0585:phyxminiproblems-problemphyxmini0585-captureate3-has-least-photonenergy}
  \lean{PhyXMiniProblems.ProblemPhyXMini0585.captureAtE3_has_least_photonEnergy}
  \uses{def:physics:phyx-mini-0585:phyxminiproblems-problemphyxmini0585-electronfinitetrapsetup, def:physics:phyx-mini-0585:phyxminiproblems-problemphyxmini0585-matchesfinitepotentialtrapscenario, def:physics:phyx-mini-0585:phyxminiproblems-problemphyxmini0585-matchesproblemreadouts, def:physics:phyx-mini-0585:phyxminiproblems-problemphyxmini0585-hasphysicalfinitetrapparameters, def:physics:phyx-mini-0585:phyxminiproblems-problemphyxmini0585-satisfieselectronentryandcapturelaws, def:physics:phyx-mini-0585:phyxminiproblems-problemphyxmini0585-isleastcapturephotonenergy}
  Since E₃ is the highest bound level, capture there removes less energy than capture at E₁ or E₂; its emitted photon therefore realizes the least possible single-photon capture energy
\end{lemma}
\begin{proof}
  The conclusion follows from the governing relations and figure readouts named in the statement of capture At E3 has least photon Energy.
\end{proof}
\begin{definition}[Answer Choice]
  \label{def:physics:phyx-mini-0585:phyxminiproblems-problemphyxmini0585-answerchoice}
  \lean{PhyXMiniProblems.ProblemPhyXMini0585.AnswerChoice}
  Labels printed beside the four energy choices
\end{definition}
\begin{proof}
  This is the declared model definition of Answer Choice.
\end{proof}
\begin{definition}[energy In Electron Volts]
  \label{def:physics:phyx-mini-0585:phyxminiproblems-problemphyxmini0585-answerchoice-energyinelectronvolts}
  \lean{PhyXMiniProblems.ProblemPhyXMini0585.AnswerChoice.energyInElectronVolts}
  \uses{def:physics:phyx-mini-0585:phyxminiproblems-problemphyxmini0585-energyinelectronvolts, def:physics:phyx-mini-0585:phyxminiproblems-problemphyxmini0585-answerchoice}
  Displayed scalar energy for each choice, in electron volts
\end{definition}
\begin{proof}
  This is the declared model definition of energy In Electron Volts.
\end{proof}
\begin{definition}[Is Correct For]
  \label{def:physics:phyx-mini-0585:phyxminiproblems-problemphyxmini0585-answerchoice-iscorrectfor}
  \lean{PhyXMiniProblems.ProblemPhyXMini0585.AnswerChoice.IsCorrectFor}
  \uses{def:physics:phyx-mini-0585:phyxminiproblems-problemphyxmini0585-energyinelectronvolts, def:physics:phyx-mini-0585:phyxminiproblems-problemphyxmini0585-electronfinitetrapsetup, def:physics:phyx-mini-0585:phyxminiproblems-problemphyxmini0585-isleastcapturephotonenergy, def:physics:phyx-mini-0585:phyxminiproblems-problemphyxmini0585-answerchoice}
  An answer choice is physically correct when its displayed value is realized by a least-energy capture photon. This remains a substantive proposition: the least photon observable is not defined from the choice table
\end{definition}
\begin{proof}
  This is the declared model definition of Is Correct For.
\end{proof}
\begin{definition}[recorded Dataset Answer]
  \label{def:physics:phyx-mini-0585:phyxminiproblems-problemphyxmini0585-recordeddatasetanswer}
  \lean{PhyXMiniProblems.ProblemPhyXMini0585.recordedDatasetAnswer}
  \uses{def:physics:phyx-mini-0585:phyxminiproblems-problemphyxmini0585-answerchoice}
  The answer label recorded in the source dataset
\end{definition}
\begin{proof}
  This is the declared model definition of recorded Dataset Answer.
\end{proof}
% --- Archon declaration topology end ---
% --- Archon physics formalization source end ---
